\documentclass[10pt,a4paper]{amsart}
\usepackage[margin=19mm]{geometry}
\usepackage{amsmath,amssymb,mathtools}
\usepackage[scr=rsfs,cal=euler]{mathalfa}
\usepackage{xfrac}
\usepackage[unicode,hidelinks,bookmarks=false]{hyperref}
\newcommand{\ie}{i.e.\ }
\newcommand{\cf}{cf.\ }
\newcommand{\resp}{respectively\ }
\newcommand{\Subsection}{\subsection}
\providecommand{\nobreakdash}{\nobreak\hbox{-}}
\setlength{\emergencystretch}{2em}
\multlinegap0pt
\usepackage{tikz}
\usetikzlibrary{cd}
\usepackage{caption}
\usepackage{amssymb}
\usepackage{float}
\newtheorem{lemma}{Lemma}
\usepackage{cleveref}
\crefname{lemma}{Lemma}{Lemmas}
\newcommand{\cb}{\mathcal{B}}
\newcommand{\cbr}{\mathcal{B}_R}
\newcommand{\co}{\mathcal{O}}
\newcommand{\ga}{\mathfrak{a}}
\newcommand{\gb}{\mathfrak{b}}
\newcommand{\gc}{\mathfrak{c}}
\title{Sarnak's Program for Erd\H{o}s Sieves. Part II: Measure Systems and Applications}
\author{Francisco Ara\'ujo}
\date{Source: arXiv:2602.24034v1 (CC BY 4.0)}
\begin{document}
\maketitle

\noindent\emph{Source and licence.} Sarnak's Program for Erd\H{o}s Sieves. Part II: Measure Systems and Applications. Version arXiv:2602.24034v1 (CC BY 4.0), distributed by arXiv under the Creative Commons Attribution 4.0 International licence, http://creativecommons.org/licenses/by/4.0/, as recorded in the arXiv licence metadata for this exact version and shown on the abstract page https://arxiv.org/abs/2602.24034v1. Full licence notice: \texttt{licenses/arxiv-2602.24034-CC-BY-4.0.txt}.\par\smallskip

\noindent\textit{Editorial scope.} Counted unit: the article's lemma base (source0.tex lines 876--879). The article's complete original proof of it is delivered with it (source0.tex lines 881--929, one \texttt{proof} environment of 49 lines). No mathematics is written, repaired or re-derived: the statement and the proof are the article's own lines, in the article's own notation, and no retained line is substituted or re-flowed.  Retained uncounted as the source-local context the proof needs: lines 578--581; lines 786--789; lines 840--843.  Nothing was omitted from any retained span, so the package carries no omission record.  No retained line contains a citation, so the wrapper carries no bibliography and no bibliography entry.  The retained lines contain the article's own diagram code, which the wrapper typesets with the standard tikz packages; no image file is delivered and the package declares no asset dependency.  Editorial formatting only, in addition: an AMS \texttt{amsart} preamble that restates the article's own declarations and macros used by the retained lines, namely the environments \texttt{lemma} on separate counters, together with the standard packages those lines need.  The retained environments print in this document as Lemma 1 in order of appearance, whereas the article prints the same items with the numbering of its own full document.  The complete native source of the article as distributed in the arXiv e-print for this exact version is bundled unchanged as \texttt{sources/195441/source0.tex}.
		\begin{lemma}
			\label{lm: if x gb in union must belong to one}
			Let $\gb$ be an ideal of $\co_K$, and $\gb_1,\dots,\gb_r$ a collection of ideals such that $(\gb,\gb_i) \neq 1$. Let $R_{\gb_i}$ be a collection of congruence classes modulo $\gb_i$ such that $R_{\gb_i} \neq \co_K$. If $$x+\gb \subset \bigcup_{i=1}^r R_{\gb_i},$$ then there is some $j$ such that $x+\gb \subset R_{\gb_j}$.
		\end{lemma}

			\begin{lemma}
				\label{lm: there are admissible representatives of R_gb^c  }
				Let $R$ be an Erdős sieve. For every $\gb \in \cb_R$, there is some $R-$admissible set $A$ such that $A+\gb = R_\gb^c$. 
			\end{lemma}

			\begin{lemma}
				\label{lm: isomorphism of dynamical systems for contraction dilation}
				Let $R$ be an Erdős sieve, and $R'$ a dilation of $R$. Then $\Omega_R = \Omega_{R'}$ and $\nu_R = \nu_{R'}.$ In particular, the identity map is an isomorphism of the dynamical systems $(\Omega_{R},S,\nu_{R})$ and $(\Omega_{R'},S,\nu_{R'})$.
			\end{lemma}

			\begin{lemma}
				\label{lm: OmegaR =omegar' implies equality of base}
				Let $R$ and $R'$ be minimal Erdős sieves. If $\Omega_{R} = \Omega_{R'}$, then $\cb_R = \cb_{R'}$. 
			\end{lemma}

			\begin{proof}
				Take any $\gb \in \cb_R$. We claim that there must be some $\gb' \in \cb_{R'}$ such that $(\gb,\gb') \neq 1$. To show this we assume that  $(\gb,\gb') = 1$ for every $\gb' \in \cb_{R'}$, and show that this implies that $\Omega_{R} \neq \Omega_{R'}$. In order to do so, we proceed similarly to how we did in the proof of  \Cref{lm: there are admissible representatives of R_gb^c  } to find some $A \in \Omega_{R'}$ that is not in $\Omega_R$.
				
				We consider $$\Delta := \{\gb'\in \cb_{R'}: \frac{|R'_{\gb'}|}{N(\gb')} \geq \frac{1}{N(\gb)}\}.$$ Since $R'$ is Erdős, $\Delta$ is finite. Using the Chinese Remainder Theorem, we can find a finite set $A$ such that $A+ \gb = \co_K$, and $A \subset (R'_{\gb'})^c$ for every $\gb' \in \Delta$. This means that $A \not \in \Omega_R$, but we now show that it is in $\Omega_{R'}$. For ideals $\gb\in \Delta$, we  know by definition of $A$ that $A \subset (R'_{\gb'})^c$. For ideals $\gb \not \in \Delta$ we procced as in in  \Cref{lm: there are admissible representatives of R_gb^c  }. For every such ideal we have that $|A|\frac{|R'_{\gb'}|}{N(\gb')}  < 1$, so $A+ R'_{\gb'}$ cannot cover $\co_K$, and so  $A \in \Omega_{R'}$. We obtain the desired contradiction, so for any $\gb \in \cbr$, there must be some $\gb' \in \cb_{R'}$ such that $(\gb,\gb') \neq 1$.
				
				We now write $V_{\gb} := \{\gb' \in \cb_{R'}: (\gb, \gb') \neq 1 \},$ which must be a non-empty set. We will show that it equals $\{\gb\}$. The first step is to show that there exists some $x \in \co_K$ such that \begin{equation}
					\label{eq: Rgb subset of union}
					R_\gb \subset x+ \bigcup_{\gb' \in V_\gb} R'_{\gb'}.
				\end{equation}
				Writing $\gc = \text{lcm}(\{\gb\}\cup V_\gb)$, this is equivalent to showing that $$R_\gb + \gc \subset x+ \bigcup_{\gb' \in V_\gb} R'_{\gb'} + \gc.$$
				
				Consider the sieve $W$ supported on $\cb_W = (\cb_{R'}\cup\{\gc\})\setminus V_\gb$ and defined by $W_\gc = \bigcup_{\gb' \in V_\gb} R'_{\gb'} + \gc$ and $W_{\gb'} = R'_{\gb'}$ for any other $\gb' \in \cb_W$. This is a dilation of $R'$, therefore \Cref{lm: isomorphism of dynamical systems for contraction dilation} implies that $\Omega_W = \Omega_{R'} = \Omega_{R}$. By \Cref{lm: there are admissible representatives of R_gb^c  }, we can find some finite $A \in \Omega_W$ such that $A+\gc = (W_\gc)^c$. Assume that for all $x\in \co_K$, we have $$R_\gb +\gc \not \subset x+ \bigcup_{\gb' \in V_\gb} R'_{\gb'} + \gc =  x + W_\gc.$$ Then we get that $(R_\gb + \gc) \cap x + (W_\gc)^c \neq \emptyset$ for all $x\in \co_K$. This implies that $$\emptyset \neq R_\gb \cap (x+A+\gc) \subset R_\gb \cap (x+A+\gb) $$ for all $x\in \co_K$, which shows that $A$ is not in $\Omega_{R}$, contradicting the fact that $\Omega_W = \Omega_{R}$. Consequently, there must be some $x\in\co_K$ such that \Cref{eq: Rgb subset of union} holds, as we wanted to show.
				
				'Visually', the proof now looks as follows. As in Section 6, we consider a bipartite graph with edges labeled $R_\gb$ or $R'_{\gb'}$, with an edge between $R_\gb$ and  $R'_{\gb'}$ if $(\gb,\gb') \neq 1$. We first claim that, writing $V_\gb$ as $\{\gb'_1,\dots,\gb'_r\}$ and taking some $\ga \in \cb_R$ different from $\gb$, this graph cannot have a subgraph that looks as follows.
				
				% https://q.uiver.app/#q=WzAsNixbMCwwLCJSX1xcZ2IiXSxbMiwwLCJSJ197XFxnYidfMX0iXSxbMiwxLCJSJ197XFxnYidfMn0iXSxbMiwzLCJSJ197XFxnYidfcn0iXSxbMiwyLCJcXHZkb3RzIl0sWzAsMiwiUl9cXGdhIl0sWzAsMSwiIiwwLHsic3R5bGUiOnsiaGVhZCI6eyJuYW1lIjoibm9uZSJ9fX1dLFsyLDAsIiIsMCx7InN0eWxlIjp7ImhlYWQiOnsibmFtZSI6Im5vbmUifX19XSxbMywwLCIiLDAseyJzdHlsZSI6eyJoZWFkIjp7Im5hbWUiOiJub25lIn19fV0sWzEsNSwiIiwwLHsic3R5bGUiOnsiaGVhZCI6eyJuYW1lIjoibm9uZSJ9fX1dXQ==
				\[\begin{tikzcd}
					{R_\gb} && {R'_{\gb'_1}} \\
					&& {R'_{\gb'_2}} \\
					{R_\ga} && \vdots \\
					&& {R'_{\gb'_r}}
					\arrow[no head, from=1-1, to=1-3]
					\arrow[no head, from=1-3, to=3-1]
					\arrow[no head, from=2-3, to=1-1]
					\arrow[no head, from=4-3, to=1-1]
				\end{tikzcd}\]
				
				This means that this graph can be written as the union of disjoint graphs of the form 
				
				% https://q.uiver.app/#q=WzAsNSxbMCwwLCJSX1xcZ2IiXSxbMiwwLCJSJ197XFxnYidfMX0iXSxbMiwxLCJSJ197XFxnYidfMn0iXSxbMiwzLCJSJ197XFxnYidfcn0iXSxbMiwyLCJcXHZkb3RzIl0sWzAsMSwiIiwwLHsic3R5bGUiOnsiaGVhZCI6eyJuYW1lIjoibm9uZSJ9fX1dLFsyLDAsIiIsMCx7InN0eWxlIjp7ImhlYWQiOnsibmFtZSI6Im5vbmUifX19XSxbMywwLCIiLDAseyJzdHlsZSI6eyJoZWFkIjp7Im5hbWUiOiJub25lIn19fV1d
				\[\begin{tikzcd}
					{R_\gb} && {R'_{\gb'_1}} \\
					&& {R'_{\gb'_2}} \\
					&& \vdots \\
					&& {R'_{\gb'_r}}
					\arrow[no head, from=1-1, to=1-3]
					\arrow[no head, from=2-3, to=1-1]
					\arrow[no head, from=4-3, to=1-1]
				\end{tikzcd}\] or the equivalent mirrored graphs (meaning that for some $R'_{\gb'}$, it will be connected to some $R_{\gb_1}, \dots,R_{\gb_l}$), and then we will use the minimality of $R$ and $R'$ to show that we must have $r = 1$.
				
				More formally, take $y \in R_\gb$. By \Cref{eq: Rgb subset of union}, there is some $x$ (not depending on $y$) such that $y-x+\gb \subset \bigcup_{\gb' \in V_\gb} R'_{\gb'}$. By \Cref{lm: if x gb in union must belong to one}, it follows that $y-x + \gb \subset R'_{\gb'_i}$ for some $i$. Consequently, we must have that $y-x+(\gb,\gb'_i) \subset R'_{\gb'_i}$. We now claim that there must be some $x_i$ such that $y-x_i+(\gb,\gb'_i) \subset R_\gb$ if $y-x \subset R'_{\gb'_i}$.
				
				Indeed, assume that $y-x-t+(\gb,\gb'_i) \not \subset R_\gb$ for all $t\in \co_K$. We know that there is some $t'_i\in \co_K$ such that $$y-x+(\gb,\gb'_i) \subset R'_{\gb'_i} \subset t'_i+\bigcup_{\ga \in \cb_R: (\ga,\gb'_i) \neq 1 }R_\ga,$$ using \Cref{eq: Rgb subset of union} applied to $R'_{\gb'_i}$. If $y-x-t'_i+(\gb,\gb'_i) \not \subset R_\gb$, then there must be some $z \in (\gb,\gb'_i)$ such that $y-x-t'_i+z+\gb \not \subset R_\gb$. Consequently, we would have, using \Cref{lm: if x gb in union must belong to one}, that  $y-x-t'_i+\gb \subset R_\ga$ for some $\ga \in \cb_R$, distinct from $\gb$. But this is impossible, since this would imply that $\co_K = \ga+\gb \subset R_\ga$. It follows that if $y-x \subset R'_{\gb'_i}$, then taking $x_i = x+t'_i$, we must have $y-x_i + (\gb,\gb'_i) \subset R_\gb$.
				
				Choosing representatives $y_1,\dots,y_{|R_\gb|}$ for $R_\gb$ in $\co_K$, let $A_i$ be the set of $y_j$ such that $y_j + \gb \subset R'_{\gb'_i}$. Using the Chinese Remainder Theorem, we can find some $x$ such that $x \equiv x_i \mod (\gb,\gb'_i)$ for all $i$ (with the $x_i$ whose existence we showed in the last paragraph). We get that $$R_\gb \subset \bigcup_{i=1}^r A_i + (\gb,\gb_i') = \bigcup_{i=1}^r A_i +x-x_i + (\gb,\gb_i') = x+\bigcup_{i=1}^r A_i -x_i+ (\gb,\gb_i')  \subset x+ R_\gb.$$ As finite subsets of $\co_K/\gb$, both $R_\gb$ and $x+R_\gb$ have the same cardinality, so $R_\gb \subset x+R_\gb$ implies equality. Consequently, we get that
				$$\bigcup_{i=1}^r A_i + (\gb,\gb_i') = R_\gb.$$ 
				
				By minimality of $R$, we conclude that there must be some $i$ such that $ (\gb,\gb_i') = \gb$, which by coprimality of the $\gb'_i$, shows that there must be a unique $\gb'$ such that $V_\gb = \{\gb'\}$ and $\gb \mid \gb'$. Using the minimality of $R'$, we can use the same argument to show that the set of those $\ga \in \cb_R$ such that $(\ga,\gb') \neq 1$ must also have only one element which is divisible by $\gb'$. This unique element must be $\gb$, and since $\gb \mid \gb'$ and $\gb' \mid \gb$, we must have an equality $\gb = \gb'$, which shows that $\cb_R = \cb_{R'}$ as we wanted to show.
			\end{proof}

\end{document}
